% -*-coding: utf-8 -*-

\defaultfont

\BiAppendixChapter{致~~~~谢}{Acknowledgements}


行文至此，落笔为终。四年的本科时光如白驹过隙，随着毕业论文的定稿，我在哈尔滨工程大学的学习生涯也即将画上一个圆满的句号。回望这四年的求学之路，有深夜面对屏幕调试代码的疲惫，更有见证系统成功运行时的喜悦。在此，我谨向所有在这个过程中给予我指导、帮助与陪伴的人，表达最深切的感激。

首先，我要向我的毕设导师刘星老师致以最诚挚的谢意。在毕业设计的全过程中，从初期结构方案的推敲研讨、理论建模的反复修改，到最终软硬件闭环联调的攻坚阶段，刘老师都倾注了大量心血。您严谨求实的治学态度、深厚的学术造诣以及对细节的精益求精，不仅让我的论文得以顺利定稿，更在潜移默化中塑造了我作为一名工科生应有的专业素养。

其次，我要特别感谢张永锐老师。是您为我推开了科技创新的大门，将我真正领入了工程实践的广阔领域。在过往的科创项目与系统设计中，您的悉心教导与鼓励，让我从面对底层硬件束手无策的初学者，一步步积累了将纸面理论转化为物理实体的能力。这份宝贵的科创启蒙之恩，我将永远铭记于心。

同时，感谢陪伴我走过四年青春岁月的朋友们，以及机电工程学院的同窗们。感谢你们在项目遇到瓶颈、系统报错时的鼎力相助，以及在生活中的彼此包容与扶持。在实验室里共同切磋技术、探讨方案的那些日子，将是我本科阶段最难忘的记忆。

我还要将最深的感恩献给我的父母。二十多年来，是你们无私的爱、理解与默默奉献，为我筑起了最坚实的后盾。无论我面临怎样的压力与抉择，你们始终是我最温暖的避风港，给予我勇往直前、追求工程梦想的底气。

最后，衷心感谢母校的悉心培养，感谢所有评阅本论文和参加答辩的老师们在百忙之中付出的辛勤劳动。路漫漫其修远兮，吾将上下而求索。在未来的道路上，我将带着这份沉甸甸的感恩，踏实进取，继续前行。
